\documentclass[11pt]{article}
\usepackage[margin=0.9in]{geometry}
\usepackage{hyperref}
\usepackage{listings}
\usepackage{xcolor}
\usepackage{booktabs}
\usepackage{enumitem}
\usepackage{longtable}
\usepackage{array}
\usepackage{amsmath}
\usepackage{ragged2e}
\usepackage{seqsplit}

% Be much more permissive about line breaking; long \texttt{} strings
% would otherwise produce overfull hboxes in the abstract and tables.
\emergencystretch=6em
\tolerance=9999
\hbadness=10000
\sloppy

% Allow line breaks inside \texttt at any underscore.
\makeatletter
\let\old@texttt\texttt
\renewcommand{\texttt}[1]{\old@texttt{\hbadness=10000 \relax #1}}
\makeatother

\lstset{
  basicstyle=\ttfamily\small,
  breaklines=true,
  columns=fullflexible,
  frame=single,
  backgroundcolor=\color{gray!10},
  showstringspaces=false
}

% Column types that wrap and ragged-right.
\newcolumntype{L}[1]{>{\RaggedRight\arraybackslash\hangindent=1em}p{#1}}
\newcommand{\status}[1]{\textbf{#1}}
\newcommand{\miss}[1]{\textbf{#1}}
\newcommand{\guess}[1]{\textit{#1}}

\title{COARDS Audit of \texttt{SST\_Orbits/} on the URI OPeNDAP Server\\
       \large grep-dap test case, Prompt \#2}
\author{Compiled for P. Cornillon}
\date{2026-05-16}

\begin{document}
\maketitle

\begin{abstract}
This document supplements the Prompt~\#1 exploration of
\url{https://sst-aqua.gso.uri.edu/opendap/} by auditing the
COARDS-essential semantic metadata of every variable in the
publicly readable \texttt{SST\_Orbits/} branch, using the
representative orbit file
\texttt{AQUA\_MODIS\_orbit\_115220\_\allowbreak 20240101T011558\_\allowbreak L2\_SST-URI\_\allowbreak 24-2.nc4}
as the sample. For each of the 44 variables across the three groups
of the file, the audit reports whether the COARDS-essential
attributes are complete; for the 23 variables judged incomplete,
the audit lists the missing attributes and a best-guess value for
each. The audit also lists, per variable, semantic metadata not
strictly required by COARDS that would significantly facilitate
use --- typically \texttt{coordinates}, \texttt{ancillary\_variables},
\texttt{cell\_methods}, \texttt{flag\_values}/\texttt{flag\_meanings},
and provenance fields. Two factual errors in the Prompt~\#1
write-up are corrected.
\end{abstract}

\tableofcontents

\section{Background and corrections to Prompt \#1}

The Prompt~\#1 document treats \texttt{SST\_Orbits/} as
``richly CF-annotated''; the deeper read carried out for this
prompt confirms that overall assessment but corrects two specific
claims in Prompt~\#1.

\begin{enumerate}
  \item Prompt~\#1 stated that the orbit DMR exposes no root-group
    (global) attributes. The DMR in fact carries 22 root-group
    attributes (\S\ref{sec:global}) including
    \texttt{title}, \texttt{Conventions="CF-1.5"}, an ACDD-style
    discovery block, and full provenance and contact
    information. The previous omission was caused by a flat-regex
    parser that only matched attribute elements inside
    \texttt{<Variable>} bodies.
  \item Prompt~\#1 stated that the Int32 \texttt{latitude} and
    \texttt{longitude} variables in the orbit file declare
    \texttt{units} but \emph{not} a \texttt{scale\_factor}, and
    that the implicit factor of $10^{-3}$ would have to be
    inferred. The DMR in fact declares \texttt{scale\_factor=0.001}
    and \texttt{add\_offset=0} on those variables. The previous
    error was the same flat-regex parsing issue: only the first
    few attributes per variable were being captured.
\end{enumerate}

This document is built on a depth-aware DMR parser
(\texttt{scripts/sst\_orbits\_dmr\_inventory.py}) and a COARDS
scorer (\texttt{scripts/sst\_orbits\_coards\_audit.py}); both are
checked into the repository and emit the JSON file used to
generate the tables below.

\section{COARDS criteria applied}\label{sec:criteria}

COARDS is the predecessor of CF; CF subsumes it. For this audit a
deliberately strict reading of COARDS is used, namely that an
``essential'' variable carries:

\begin{description}[leftmargin=2em,topsep=2pt,style=nextline]
  \item[\texttt{units}.] A UDUNITS-compatible string. For
    latitudes one of \{\texttt{degrees\_north},
    \texttt{degree\_north}, \texttt{degree\_N},
    \texttt{degrees\_N}, \texttt{degreeN}, \texttt{degreesN}\}; for
    longitudes the corresponding east set. For absolute-time
    variables the string must be of the form
    \texttt{<unit> since <reference>}.
  \item[\texttt{long\_name}.] A human-readable description.
    COARDS recommends rather than requires it, but it is the
    single most useful attribute for an outside reader; the audit
    treats it as required.
  \item[\texttt{\_FillValue}.] For every non-scalar, non-string
    variable that may carry missing data.
  \item[\texttt{scale\_factor} \& \texttt{add\_offset}.] If
    either is set, the other must also be set; this rules out
    silently dropped scaling.
\end{description}

The audit additionally treats the units string \texttt{"C"} as a
COARDS failure --- it is not in UDUNITS (where \texttt{C} means
\emph{coulomb}); the COARDS-compatible string is
\texttt{degree\_Celsius} (or \texttt{K}). \texttt{"C/km"} is
flagged for the same reason.

\section{Global (root-group) attributes}\label{sec:global}

The orbit file's root group exposes 22 attributes. They fall into
three groups:

\sloppy
\begin{description}[leftmargin=2em,topsep=2pt,style=nextline]
  \item[Conventions and identity.]
    \texttt{title="Conditioned 1km (L2) MODIS AQUA SST"},
    \texttt{Conventions="CF-1.5"},
    \texttt{Metadata\_Conventions=}\linebreak[1]\texttt{"Unidata Dataset Discovery v1.0"},
    \texttt{standard\_name\_vocabulary="CF-1.5"},
    \texttt{uuid="03bf4c3e-7598-4264-a863-dba03d2a6d46"},
    \texttt{cdm\_data\_type="Grid"}.
  \item[Provenance.] \texttt{date\_created},
    \texttt{source="Satellite observation"},
    \texttt{history="\{17-May-2025 09:20:41 :
    Fix\_MODIS\_Mask version1:$\langle$uuid$\rangle$\}"},
    \texttt{Summary} (a free-text paragraph stating the file
    is derived from L2 MODIS Aqua granules from NASA's
    OceanColor web site),
    \texttt{contributor\_name="NASA"},
    \texttt{contributor\_role="Generated SST fields from
    MODIS data."}
  \item[Contacts and licence.] \texttt{creator\_name="Peter
    Cornillon"}, \texttt{creator\_url},
    \texttt{creator\_email},
    \texttt{institution="University of Rhode Island - Graduate
    School of Oceanography"},
    \texttt{project="SST Fronts"},
    \texttt{acknowledgement}, \texttt{license},
    \texttt{publisher\_name="SST Fronts"},
    \texttt{publisher\_url}, \texttt{publisher\_institution}.
\end{description}
\fussy

\paragraph{COARDS status.} The minimum COARDS / NUG global
attribute set is \{\texttt{Conventions}, \texttt{title},
\texttt{history}, \texttt{source}, \texttt{references},
\texttt{comment}, \texttt{institution}\}. Of these, five
(\texttt{Conventions}, \texttt{title}, \texttt{history},
\texttt{source}, \texttt{institution}) are present and one
(\texttt{references}) is absent. The \texttt{comment} attribute is
not present under that name; \texttt{Summary} fills the same role
but a CF/COARDS-aware reader would not find it under the canonical
key. \textbf{Best-guess additions:}

\begin{itemize}
  \item \texttt{Conventions = "CF-1.8"}~--- the file already
    claims CF-1.5; bumping to a recent version is harmless and
    forces a re-validation that would catch the \texttt{units="C"}
    issues described below.
  \item \texttt{references = "Cornillon, P., et al., SST Fronts
    project,
    \url{http://www.sstfronts.org}; MODIS Aqua L2 SST,
    NASA OceanColor;
    \url{https://oceancolor.gsfc.nasa.gov/}"}.
  \item \texttt{comment = $\langle$ the existing Summary text $\rangle$};
    that is, alias the \texttt{Summary} attribute under the
    standard \texttt{comment} key.
  \item \texttt{platform = "Aqua"}, \texttt{instrument =
    "MODIS"}, \texttt{processing\_level = "L2"},
    \texttt{product\_version = "URI 24-2"} --- these are not
    strictly COARDS but every modern catalog scraper (ACDD,
    GCMD) expects them.
  \item \texttt{geospatial\_lat\_min/lat\_max/lon\_min/lon\_max}
    and \texttt{time\_coverage\_start/end} (filled from the
    orbit's actual extent) --- ACDD, not COARDS, but essential
    for any harvester.
  \item \texttt{keywords = "Earth Science > Oceans > Ocean
    Temperature > Sea Surface Temperature; Earth Science >
    Oceans > Ocean Temperature > SST Fronts; Earth Science >
    Oceans > Ocean Temperature > SST Gradients"} --- GCMD
    keywords for discoverability.
\end{itemize}

\section{Per-variable assessment}

Three tables follow, one per group, listing every variable, its
COARDS status, the COARDS-essential attributes that are missing
(with a best-guess value for each), and a list of non-COARDS
attributes whose addition would significantly facilitate use.

In each table the second column reads ``Complete'' if all of
\texttt{units}, \texttt{long\_name}, \texttt{\_FillValue}, and
\texttt{scale\_factor}/\texttt{add\_offset} (where applicable) are
present and the \texttt{units} string is COARDS-compatible.
Otherwise it lists the missing items in the form
\miss{attr} $\to$ \guess{best-guess value}.

\subsection{Root group (\texttt{/})}

\footnotesize
\begin{longtable}{@{}L{3.6cm} L{5.9cm} L{6.7cm}@{}}
\toprule
\textbf{Variable} & \textbf{COARDS: status / missing $\to$ guess} & \textbf{Recommended non-COARDS additions} \\
\midrule
\endfirsthead
\toprule
\textbf{Variable} & \textbf{COARDS: status / missing $\to$ guess} & \textbf{Recommended non-COARDS additions} \\
\midrule
\endhead

\texttt{SST\_In} \newline Int16 (ny,nx)
& \status{Incomplete.} \miss{units} $\to$ \guess{\texttt{degree\_Celsius}} (the file declares \texttt{"C"}, which is UDUNITS \emph{coulomb}).
& \texttt{coordinates="latitude longitude DateTime"};
  \texttt{ancillary\_variables="qual\_sst refined\_mask"};
  \texttt{cell\_methods="area: point"};
  \texttt{source="MODIS Aqua L2 SST, NASA OceanColor"};
  \texttt{coverage\_content\_type="physicalMeasurement"}. \\

\texttt{regridded\_sst} \newline Int16 (ny,nx)
& \status{Incomplete.} \miss{units} $\to$ \guess{\texttt{degree\_Celsius}}.
& \texttt{coordinates=}\linebreak[1]\texttt{"regridded\_latitude regridded\_longitude DateTime"};
  \texttt{ancillary\_variables="qual\_sst refined\_mask"};
  \texttt{cell\_methods="area: mean"};
  \texttt{comment="SST after masking high-gradient pixels and regridding."} \\

\texttt{eastward\_gradient} \newline Int32 (ny,nx)
& \status{Incomplete.} \miss{units} $\to$ \guess{\texttt{degree\_Celsius km-1}} (\texttt{"C/km"} is not UDUNITS).
& \texttt{coordinates="latitude longitude DateTime"};
  \texttt{ancillary\_variables="refined\_mask qual\_sst"};
  \texttt{cell\_methods="ny: derivative nx: derivative (interval: 1 pixel)"};
  \texttt{comment="East component of the SST gradient computed from regridded\_sst."} \\

\texttt{northward\_gradient} \newline Int32 (ny,nx)
& \status{Incomplete.} \miss{units} $\to$ \guess{\texttt{degree\_Celsius km-1}}.
& Same as \texttt{eastward\_gradient} with \texttt{comment} adjusted to ``North component\dots'' \\

\texttt{refined\_mask} \newline Int8 (ny,nx)
& \status{Incomplete.} \miss{units} $\to$ \guess{\texttt{"1"}} (dimensionless flag).
& Replace malformed \texttt{flag\_masks=0} (a single scalar of 0 does not match CF \texttt{flag\_masks} semantics) with
  \texttt{flag\_values="0 1"} and
  \texttt{flag\_meanings="good bad"}. Add
  \texttt{coordinates="latitude longitude DateTime"} and
  \texttt{standard\_name="sea\_surface\_temperature status\_flag"}. \\

\texttt{qual\_sst} \newline Int8 (ny,nx)
& \status{Incomplete.} \miss{units} $\to$ \guess{\texttt{"1"}}.
& Replace \texttt{flag\_masks=0} with
  \texttt{flag\_values="0 1 2 3 4"} and
  \texttt{flag\_meanings="best good questionable bad notprocessed"}.
  Rename \texttt{standard\_name} from
  \texttt{"sea\_surface\_temperature"} (which is incorrect for a
  flag) to
  \texttt{"sea\_surface\_temperature status\_flag"}. Add
  \texttt{coordinates="latitude longitude DateTime"}. \\

\texttt{latitude} \newline Int32 (ny,nx)
& \status{Complete.} \texttt{units=degrees\_north},
  \texttt{scale\_factor=0.001}, \texttt{add\_offset=0},
  \texttt{\_FillValue=-2147483647},
  \texttt{long\_name="latitude"}, \texttt{valid\_min/max} set.
& \texttt{axis="Y"}; could add
  \texttt{comment="Per-pixel geodetic latitude on WGS-84."} \\

\texttt{longitude} \newline Int32 (ny,nx)
& \status{Complete.} \texttt{units=degrees\_east},
  \texttt{scale\_factor=0.001}, \texttt{add\_offset=0},
  \texttt{\_FillValue} present, \texttt{long\_name="Longitude"},
  \texttt{valid\_min/max=$\pm$720000}. The valid range of
  $\pm 720^\circ$ permits unwrapped longitudes along an orbit ---
  unusual but coherent.
& \texttt{axis="X"};
  \texttt{comment="Unwrapped longitude; may exceed $\pm 180^\circ$ within a single orbit."} \\

\texttt{regridded\_latitude} \newline Int32 (ny,nx)
& \status{Complete.}
& \texttt{axis="Y"}; \texttt{comment="Latitude after the regridding step that produced \texttt{regridded\_sst}."} \\

\texttt{regridded\_longitude} \newline Int32 (ny,nx)
& \status{Complete.}
& \texttt{axis="X"}; \texttt{comment} analogous to above. \\

\texttt{nadir\_latitude} \newline Float32 (ny)
& \status{Complete.}
& \texttt{axis="Y"};
  \texttt{coordinates="DateTime"};
  \texttt{comment="Sub-spacecraft (nadir) latitude per scan line."} \\

\texttt{nadir\_longitude} \newline Float32 (ny)
& \status{Complete.}
& \texttt{axis="X"}; \texttt{comment} analogous. \\

\texttt{left\_swath\_edge\_}\linebreak[1]\texttt{trackline\_latitude} \newline Float32 (ny)
& \status{Complete.}
& \texttt{axis="Y"};
  \texttt{comment="Geodetic latitude of the left (port-side) swath edge per scan."} \\

\texttt{right\_swath\_edge\_}\linebreak[1]\texttt{trackline\_latitude} \newline Float32 (ny)
& \status{Complete.}
& \texttt{axis="Y"}; analogous comment. \\

\texttt{left\_swath\_edge\_}\linebreak[1]\texttt{trackline\_longitude} \newline Float32 (ny)
& \status{Complete.}
& \texttt{axis="X"}; analogous comment. \\

\texttt{right\_swath\_edge\_}\linebreak[1]\texttt{trackline\_longitude} \newline Float32 (ny)
& \status{Complete.}
& \texttt{axis="X"}; analogous comment. \\

\texttt{time\_from\_start\_orbit} \newline Float32 (ny)
& \status{Complete} as a duration (\texttt{units="seconds"} is acceptable for elapsed time).
& \texttt{standard\_name="time"} is technically wrong (this is an elapsed-time, not an absolute time); a closer fit would be \texttt{standard\_name="forecast\_period"} (CF lacks a clean ``time-since-orbit-start'' name) and a \texttt{comment="Seconds since the scalar DateTime (orbit start, 1970-01-01 epoch seconds)."}.\newline
  Also: \texttt{axis="T"} is debatable on a 1-D auxiliary
  coordinate but is consistent with CF practice. \\

\texttt{DateTime} \newline Float64 scalar
& \status{Incomplete.} \miss{units} $\to$ \guess{\texttt{seconds since 1970-01-01 00:00:00 UTC}} (current \texttt{units="seconds"} lacks the required reference; the reference is stated only in the \texttt{long\_name}).
& \texttt{calendar="gregorian"}; \texttt{axis="T"};
  \texttt{long\_name="Orbit start time"} (clearer than the current
  ``time since 1970-01-01 00:00:00.0''). \\

\texttt{region\_start} \newline Int32 (i)
& \status{Incomplete.} \miss{units} $\to$ \guess{\texttt{"1"}} (dimensionless index);
  \miss{\_FillValue} $\to$ \guess{\texttt{-2147483647}} (matches the
  Int32 convention used elsewhere).
& \texttt{long\_name="Start \texttt{ny}-index of each of the 5 along-orbit regions."};
  \texttt{comment="Pairs with region\_end to delimit 5 along-orbit subregions; useful for restricting processing."} \\

\texttt{region\_end} \newline Int32 (i)
& \status{Incomplete.} Same two items as \texttt{region\_start}.
& Analogous to \texttt{region\_start}. \\
\bottomrule
\end{longtable}

\subsection{Group \texttt{/Regrid\_to\_L2eqa}}

This group has the additional caveat that the \texttt{nyAMSR\_E}
and \texttt{nxAMSR\_E} dimensions are both 1 in the sample file:
AMSR-E stopped producing data in October 2011, so the AMSR-E
variables in a 2024 file are $1\times 1$ placeholders, not
populated arrays. The schema persists past the instrument's
lifetime. A reader without that context will see an apparently
intentional $1\times 1$ field and may treat it as authoritative;
adding a \texttt{comment} attribute is the cheapest fix.

\footnotesize
\begin{longtable}{@{}L{3.8cm} L{5.8cm} L{6.6cm}@{}}
\toprule
\textbf{Variable} & \textbf{COARDS: status / missing $\to$ guess} & \textbf{Recommended non-COARDS additions} \\
\midrule
\endfirsthead
\toprule
\textbf{Variable} & \textbf{COARDS: status / missing $\to$ guess} & \textbf{Recommended non-COARDS additions} \\
\midrule
\endhead

\texttt{L2eqa\_MODIS\_SST} \newline Int16 (nyL2eqa,nxL2eqa)
& \status{Incomplete.} \miss{units} $\to$ \guess{\texttt{degree\_Celsius}}.
& \texttt{coordinates="L2eqaLat L2eqaLon DateTime"};
  \texttt{cell\_methods="area: mean"};
  \texttt{ancillary\_variables="L2eqa\_MODIS\_std\_SST L2eqa\_MODIS\_num\_SST"};
  \texttt{comment="MODIS SST averaged into an L2 equal-area cell."} \\

\texttt{L2eqa\_MODIS\_std\_SST} \newline Int16 (nyL2eqa,nxL2eqa)
& \status{Incomplete.} \miss{units} $\to$ \guess{\texttt{degree\_Celsius}}.
& \texttt{coordinates="L2eqaLat L2eqaLon DateTime"};
  \texttt{cell\_methods="area: standard\_deviation"};
  \texttt{comment="Within-cell standard deviation of MODIS SST contributions."} \\

\texttt{L2eqa\_MODIS\_num\_SST} \newline Int16 (nyL2eqa,nxL2eqa)
& \status{Incomplete.} \miss{units} $\to$ \guess{\texttt{"1"}} (the
  declared \texttt{units="C"} is wrong; this is a pixel count, not
  a temperature).
& \texttt{coordinates="L2eqaLat L2eqaLon DateTime"};
  \texttt{cell\_methods="area: count"};
  \texttt{long\_name="Number of valid MODIS SST contributions in the L2eqa cell"} (clearer than current). \\

\texttt{L2eqa\_AMSR\_E\_SST} \newline Int16 (nyL2eqa,nxL2eqa)
& \status{Incomplete.} \miss{units} $\to$ \guess{\texttt{degree\_Celsius}}.
& \texttt{coordinates="L2eqaLat L2eqaLon DateTime"};
  \texttt{source="AMSR-E aboard NASA Aqua (operational 2002-06 to 2011-10)"};
  \texttt{comment="In files past 2011-10 the underlying instrument is retired; field is a 1x1 placeholder."} \\

\texttt{L2eqa\_AMSR\_E\_wind\_speed} \newline Int16 (nyL2eqa,nxL2eqa)
& \status{Incomplete.} \miss{units} $\to$ \guess{\texttt{m s-1}} (the declared
  \texttt{units="mm"} is almost certainly wrong; the
  packing \texttt{scale\_factor=0.2}, \texttt{add\_offset=25.4}
  yields a 0--50.8 range, consistent with wind speed in m s$^{-1}$
  and matching the sibling \texttt{AMSR\_E\_wind\_speed} which is
  labelled \texttt{m s-1} natively).
& \texttt{coordinates="L2eqaLat L2eqaLon DateTime"};
  \texttt{source="AMSR-E aboard NASA Aqua (operational 2002-06 to 2011-10)"};
  \texttt{comment} on AMSR-E retirement, as above. \\

\texttt{L2eqa\_AMSR\_E\_water\_vapor} \newline Int16 (nyL2eqa,nxL2eqa)
& \status{Complete.}
& \texttt{coordinates="L2eqaLat L2eqaLon DateTime"};
  retirement comment. \\

\texttt{L2eqa\_AMSR\_E\_}\linebreak[1]\texttt{cloud\_liquid\_water} \newline Int16 (nyL2eqa,nxL2eqa)
& \status{Complete.}
& As above. \\

\texttt{L2eqaLat} \newline Int32 (ny,nx) [in group]
& \status{Complete.} \texttt{units=degrees\_north},
  \texttt{scale\_factor=0.001}.
& \texttt{axis="Y"};
  \texttt{comment="Latitude of L2 equal-area cell centroid."} \\

\texttt{L2eqaLon} \newline Int32 (nyL2eqa,nxL2eqa)
& \status{Complete.} \texttt{units=degrees\_east},
  \texttt{scale\_factor=0.001}.
& \texttt{axis="X"}; analogous comment. \\

\texttt{AMSR\_E\_SST} \newline Int16 (nyAMSR\_E,nxAMSR\_E)
& \status{Incomplete.} \miss{units} $\to$ \guess{\texttt{degree\_Celsius}}.
& \texttt{coordinates="AMSR\_E\_lat AMSR\_E\_lon DateTime"};
  \texttt{source="AMSR-E aboard NASA Aqua"};
  retirement comment. \\

\texttt{AMSR\_E\_wind\_speed} \newline Int16 (nyAMSR\_E,nxAMSR\_E)
& \status{Complete.} \texttt{units="m s-1"} is UDUNITS-compatible.
& \texttt{coordinates="AMSR\_E\_lat AMSR\_E\_lon DateTime"};
  retirement comment. \\

\texttt{AMSR\_E\_water\_vapor} \newline Int16 (nyAMSR\_E,nxAMSR\_E)
& \status{Complete.}
& \texttt{coordinates="AMSR\_E\_lat AMSR\_E\_lon DateTime"};
  retirement comment;
  \texttt{long\_name="AMSR-E columnar water vapor"} (more specific than current). \\

\texttt{AMSR\_E\_cloud\_liquid\_water} \newline Int16 (nyAMSR\_E,nxAMSR\_E)
& \status{Complete.}
& As above. \\

\texttt{AMSR\_E\_lat} \newline Int32 (nyAMSR\_E,nxAMSR\_E)
& \status{Complete.}
& \texttt{axis="Y"}; retirement comment. \\

\texttt{AMSR\_E\_lon} \newline Int32 (nyAMSR\_E,nxAMSR\_E)
& \status{Complete.}
& \texttt{axis="X"}; retirement comment. \\

\texttt{MODIS\_SST\_on\_AMSR\_E\_grid} \newline Int16 (nyAMSR\_E,nxAMSR\_E)
& \status{Incomplete.} \miss{units} $\to$ \guess{\texttt{degree\_Celsius}}.
& \texttt{coordinates="AMSR\_E\_lat AMSR\_E\_lon DateTime"};
  \texttt{cell\_methods="area: mean"};
  \texttt{comment="MODIS L2eqa SST resampled onto the AMSR-E footprint."} \\

\texttt{convolved\_MODIS\_}\linebreak[1]\texttt{on\_AMSR\_E\_grid} \newline Int16 (nyAMSR\_E,nxAMSR\_E)
& \status{Incomplete.} \miss{units} $\to$ \guess{\texttt{degree\_Celsius}}.
& \texttt{coordinates="AMSR\_E\_lat AMSR\_E\_lon DateTime"};
  \texttt{cell\_methods="area: convolution"};
  \texttt{comment="MODIS L2eqa SST convolved with the AMSR-E antenna pattern; the apples-to-apples comparison field for AMSR\_E\_SST."} \\
\bottomrule
\end{longtable}

\subsection{Group \texttt{/contributing\_granules}}

This group is metadata about the L2 input granules stitched into the
orbit. It is the weakest annotated group in the file; six of seven
variables fail the strict audit.

\begin{longtable}{@{}L{3.0cm} L{6.4cm} L{6.8cm}@{}}
\toprule
\textbf{Variable} & \textbf{COARDS: status / missing $\to$ guess} & \textbf{Recommended non-COARDS additions} \\
\midrule
\endfirsthead
\toprule
\textbf{Variable} & \textbf{COARDS: status / missing $\to$ guess} & \textbf{Recommended non-COARDS additions} \\
\midrule
\endhead

\texttt{filenames} \newline String (name\_length,nGranules)
& \status{Complete.} (String variables need no \texttt{units} or \texttt{\_FillValue}.)
& \texttt{long\_name=}\linebreak[1]\texttt{"Source MODIS L2 granule filenames"};
  \texttt{comment="From NASA OceanData"}\linebreak[1] \texttt{(oceandata.sci.gsfc.nasa.gov)}. \\

\texttt{start\_time} \newline Float64 (nGranules)
& \status{Incomplete.} \miss{units} $\to$ \guess{\texttt{seconds since 1970-01-01 00:00:00 UTC}} (current \texttt{units="seconds"} lacks the reference; the reference is hidden in the \texttt{long\_name});
  \miss{\_FillValue} $\to$ \guess{\texttt{NaN}}.
& \texttt{calendar="gregorian"};
  \texttt{long\_name="Granule start time"}. \\

\texttt{end\_time} \newline Float64 (nGranules)
& \status{Incomplete.} Same two items as \texttt{start\_time}.
& \texttt{calendar="gregorian"};
  \texttt{long\_name="Granule end time"}. \\

\texttt{granule\_start\_index} \newline Float64 \linebreak[1](max\_indices,\linebreak[1]nGranules)
& \status{Incomplete.} \miss{units} $\to$ \guess{\texttt{"1"}} (a dimensionless index).
  Note: \texttt{\_FillValue=NaN} \emph{is} present.
& \texttt{long\_name="First scan-line index within the source granule contributing to this orbit."} \\

\texttt{granule\_end\_index} \newline Float64 \linebreak[1](max\_indices,\linebreak[1]nGranules)
& \status{Incomplete.} \miss{units} $\to$ \guess{\texttt{"1"}}.
& \texttt{long\_name="Last scan-line index within the source granule contributing to this orbit."} \\

\texttt{orbit\_start\_index} \newline Float64 \linebreak[1](max\_indices,\linebreak[1]nGranules)
& \status{Incomplete.} \miss{units} $\to$ \guess{\texttt{"1"}}.
& \texttt{long\_name="Index in the orbit's ny axis where this granule's first scan lands."} \\

\texttt{orbit\_end\_index} \newline Float64 \linebreak[1](max\_indices,\linebreak[1]nGranules)
& \status{Incomplete.} \miss{units} $\to$ \guess{\texttt{"1"}}.
& \texttt{long\_name="Index in the orbit's ny axis where this granule's last scan lands."} \\
\bottomrule
\end{longtable}

\section{Summary}\label{sec:summary}

The orbit file \emph{looks} well annotated, and on the most
important variables it largely is: the SST fields, the gradient
fields, the lat/lon arrays, the scalar time, and the AMSR-E
co-located fields all carry \texttt{long\_name},
\texttt{standard\_name}, \texttt{\_FillValue}, \texttt{scale\_factor},
\texttt{add\_offset}, and \texttt{valid\_min/valid\_max}. The
COARDS audit nonetheless judges 23 of 44 variables incomplete,
driven by three recurring issues:

\begin{enumerate}
  \item \textbf{Units string \texttt{"C"} / \texttt{"C/km"}.} Both
    are non-UDUNITS. Substituting \texttt{degree\_Celsius} and
    \texttt{degree\_Celsius km-1} on the affected 15 variables
    would close most of the gap.
  \item \textbf{Time variables without reference.}
    \texttt{DateTime}, \texttt{start\_time},
    \texttt{end\_time} all carry \texttt{units="seconds"} with the
    epoch stated in the \texttt{long\_name} rather than the
    \texttt{units} string. A COARDS-aware client cannot infer the
    epoch from the units alone, so attempts to convert these to
    Python \texttt{datetime} or \texttt{numpy.datetime64} via
    standard tools may fail or silently produce wrong values.
    Replacing the units string with
    \texttt{"seconds since 1970-01-01 00:00:00 UTC"} (and adding
    \texttt{calendar="gregorian"}) is a near-zero-cost fix.
  \item \textbf{Dimensionless quantities with no units.} The
    \texttt{region\_*}, \texttt{*\_index}, and flag arrays
    declare no units. COARDS, and CF in particular, expects
    \texttt{units="1"} for dimensionless quantities; this is
    not optional.
\end{enumerate}

Independent of COARDS, the orbit file would be substantially
easier to read and re-use with the following additions, listed in
the per-variable tables:

\begin{itemize}
  \item \texttt{coordinates="latitude longitude DateTime"} on
    every 2-D data variable in the root group (CF-1.x links
    auxiliary coordinates to data variables this way).
  \item \texttt{ancillary\_variables="qual\_sst refined\_mask"}
    on every SST and gradient field.
  \item \texttt{cell\_methods} on derived quantities: gradients
    are differences, the regridded SST is a mean, the L2eqa fields
    are means with companion count/std, etc.
  \item Proper \texttt{flag\_values}/\texttt{flag\_meanings} on
    \texttt{refined\_mask} and \texttt{qual\_sst}. The current
    \texttt{flag\_masks=0} declaration is malformed and conveys
    no useful information.
  \item Instrument provenance: the AMSR-E retirement caveat is
    only knowable to a reader who already knows when AMSR-E
    stopped; a one-line \texttt{comment} on each AMSR-E variable
    would resolve that.
  \item Discovery attributes (\texttt{platform},
    \texttt{instrument}, \texttt{processing\_level},
    \texttt{product\_version},
    \texttt{geospatial\_*}, \texttt{time\_coverage\_*},
    \texttt{keywords}) at the root: standard ACDD additions that
    any catalog scraper expects.
\end{itemize}

\section{Caveats}

\begin{itemize}
  \item The audit was carried out on a single orbit file (the
    first orbit of 2024) chosen as representative; the
    \texttt{date\_created} of this file is 2025-05-17. Earlier
    reprocessings on the same server may carry slightly
    different attribute sets.
  \item This audit considers only the metadata. It does not
    re-check the values themselves except to verify the three
    cases where the declared units appear inconsistent with the
    packed values (\texttt{L2eqa\_AMSR\_E\_wind\_speed=mm},
    \texttt{L2eqa\_MODIS\_num\_SST=C}). Other ``Complete''
    judgements in the tables do not certify that the values are
    correct.
  \item COARDS is, in 2026, a small standard; CF supersedes it
    almost everywhere. Treating COARDS as the bar rather than CF
    is generous; an equally-strict CF-1.8 audit would flag
    everything COARDS flags here plus the missing
    \texttt{coordinates}, \texttt{ancillary\_variables},
    \texttt{cell\_methods}, and flag attributes listed in the
    rightmost columns.
\end{itemize}

\end{document}
